\section*{Chapitre \ref{ch:ml:feature-engineering-selection-and-importance} --- Construction, sélection et importance des variables}

\begin{solution}{exo:ml:feature-engineering-selection-and-importance:1}
$10^2/((2\times0.6 - 1)\times200) = 2.5$ fausses sélections attendues à 0.6~; $100/(0.8\times200) = 0.625$ à 0.9.
\end{solution}

\begin{solution}{exo:ml:feature-engineering-selection-and-importance:2}
En moyennant les contributions marginales sur les six ordres possibles, le joueur~1 contribue $1, 1, 2, 2, 1, 2$, donc $\phi_1 = 1.5$~;
le joueur~2 de même, $\phi_2 = 1.5$~; le joueur~3 n'ajoute rien dans aucun ordre, $\phi_3 = 0$. Ces valeurs somment à $v(P) = 3$.
\end{solution}

\begin{solution}{exo:ml:feature-engineering-selection-and-importance:3}
0.09 point. Cela ne montre pas que la variable est sans importance~: une autre variable peut la remplacer (l'\omterm{def:ml:feature-engineering-selection-and-importance:substitution}{effet de
substitution})~; seule l'importance de son groupe répond à cette question.
\end{solution}

\begin{solution}{exo:ml:feature-engineering-selection-and-importance:4}
Le modèle utilise les deux copies. Permutées ensemble, elles lui fournissent un bruit de la variance du signal qu'elles portaient,
si bien que ses prédictions deviennent un bruit de même ampleur, et une prédiction bruitée fait moins bien que la moyenne~: le coefficient de détermination tombe
à $2.67 - 2.72 = -0.05~\%$.
\end{solution}

\begin{solution}{exo:ml:feature-engineering-selection-and-importance:5}
La MDI additionne la réduction de la perte d'entraînement de chaque coupure, et sur les données d'entraînement une coupure sur du bruit
réduit bien la perte (elle ajuste le bruit). L'\omterm{def:ml:feature-engineering-selection-and-importance:mdi}{importance par permutation} se mesure sur des lignes mises à part, où ces coupures n'aident en rien.
\end{solution}

\begin{solution}{exo:ml:feature-engineering-selection-and-importance:6}
L'élagage variable par variable retire les copies corrélées l'une après l'autre~: chacune paraît sans importance tant que l'autre subsiste, et
le second retrait coûte alors ce qu'aucune des deux ne valait (1.05 point ici). Il faut élaguer des grappes, non des variables. Si les
importances avaient été calculées sur les données de test, la sélection constitue aussi une fuite (chapitre~3).
\end{solution}

\begin{solution}{exo:ml:feature-engineering-selection-and-importance:7}
\texttt{heavy\_tail(df=5.0)}~: brut 1.36~\%, winsorisé 1.49~\%, en rangs 1.57~\%, vrai 1.79~\%. Avec cinq degrés de liberté
les valeurs extrêmes sont rares et modérées, elles ne fixent donc plus la pente~; les transformations comptent beaucoup moins qu'avec 1.5.
\end{solution}

\begin{solution}{exo:ml:feature-engineering-selection-and-importance:8}
Si la prédiction ne dépend que de $x_1 + x_{1b}$ et que les deux variables sont identiques, la fonction valeur est symétrique en elles, et
l'axiome de symétrie donne des \omterm{def:ml:feature-engineering-selection-and-importance:shapley}{valeurs de Shapley} égales. Prenons $p = \gamma(x_1 + x_{1b}) = 2\gamma x_1$ avec $\Var(x_1) = \sigma^2$
et des résidus non corrélés aux variables. Permuter $x_1$ seule remplace $p$ par $\gamma(x_{1,\pi} + x_1)$ et ajoute
$\E[\gamma^2(x_1 - x_{1,\pi})^2] = 2\gamma^2\sigma^2$ à l'erreur quadratique moyenne~; de même pour $x_{1b}$~; permuter les deux avec
une même permutation ajoute $\E[4\gamma^2(x_1 - x_{1,\pi})^2] = 8\gamma^2\sigma^2$, soit le double de la somme des effets isolés.
\end{solution}

\begin{solution}{pb:ml:feature-engineering-selection-and-importance:1}
\begin{enumerate}
\item 0.13~\%, 1.27~\% et 1.96~\%.
\item Winsoriser à 1~\% laisse de grandes valeurs dans une distribution à queue aussi épaisse~; la transformation en rangs supprime entièrement l'échelle et
reproduit l'effet de saturation.
\item Rangs transversaux par date~: robustes aux queues et aux valeurs aberrantes, comparables d'une date à l'autre.
\item 2.67~\% hors échantillon~; la valeur vraie 4.01~\%.
\item Le gain, la permutation et la suppression de colonne classent $x_2$ en tête et $x_1$ en quatrième~; SHAP classe $x_2$, $x_1$, la copie, puis $x_3$ et
$x_4$.
\item SHAP explique le modèle ajusté, et le modèle utilise la copie.
\item 5.4~\% du total (contre 9.5~\% pour $x_1$)~: le gain se mesure sur les données d'entraînement, où les coupures sur le bruit réduisent la perte.
\item L'importance par suppression de colonne (réajustement sans la variable), et l'\omterm{def:ml:feature-engineering-selection-and-importance:mdi}{importance par permutation} sur les données mises à part seulement en partie.
\item 0.76 et 0.41 point pris isolément, 2.72 ensemble.
\item 0.09 et $-0.02$ isolément, 1.05 ensemble.
\item $\{x_1, x_{1b}\}$ à une distance de 0.5 (corrélation 0.95), toutes les autres variables seules.
\item On supprime une copie, puis l'autre, en perdant 1.05 point sur 2.67.
\item Le lasso retient $x_1$ (6.7 variables sélectionnées par demi-échantillon~; borne 2.3)~; le boosting retient $x_1$ et $x_2$ (borne 1.8)~; aucune
variable de bruit dans aucun des deux.
\item Leur effet ne passe que par le produit~: aucun criblage linéaire ne le voit, et la première coupure d'un arbre glouton sur l'un ou l'autre
ne gagne rien.
\item \emph{Résultat nommé}~: l'\omterm{def:ml:feature-engineering-selection-and-importance:mdi}{importance par permutation} et par suppression de colonne classe les signaux $x_2, x_3, x_4, x_1$ (suppression de colonne~: $x_2,
x_4, x_3, x_1$) et la copie bien plus bas~; SHAP place la copie au troisième rang~; le taux de fausses sélections de la \omterm{def:ml:feature-engineering-selection-and-importance:selection}{sélection par stabilité} à 0.6 est
nul pour les deux sélecteurs (bornes de 2.3 et 1.8 variables fausses), et les deux manquent l'interaction.
\item Un gérant de portefeuille~: importance par suppression de grappes (ce qu'exigent les données)~; un comité des risques~: SHAP pour les
décisions individuelles (ce que le modèle a fait).
\item Ajouter des produits candidats, ou utiliser un sélecteur capable de couper sur des paires (boosting avec davantage d'arbres, ou contraintes
d'interaction levées), et donner à la paire un petit effet principal à détecter.
\item Les importances doivent être calculées sur des fenêtres glissantes~; l'importance d'une variable sur tout l'échantillon moyenne des régimes.
\item Importances par grappes mises à part, avec la méthode nommée~; les copies et leur grappe~; les variables qu'aucune méthode
classe au-dessus du bruit~; la procédure de sélection et sa borne de fausses sélections.
\item La dépendance d'un modèle ajusté, ou des données, à une variable, au sens d'une seule notion qu'il faut nommer.
\end{enumerate}
\end{solution}

\begin{solution}{iq:ml:feature-engineering-selection-and-importance:1}
L'importance par impureté additionne les réductions de la perte d'entraînement des coupures sur une variable~: elle favorise les variables qui offrent beaucoup de points de
coupure et vers le bruit ajusté. L'\omterm{def:ml:feature-engineering-selection-and-importance:mdi}{importance par permutation} mesure la chute du score hors échantillon lorsque la variable est
mélangée. Il faut se fier à la mesure hors échantillon, regroupée par corrélation.
\iqlookfor{entraînement contre hors échantillon, et les biais connus de la MDI.}
\end{solution}

\begin{solution}{iq:ml:feature-engineering-selection-and-importance:2}
On les mesure comme un groupe (permuter ou supprimer les deux)~: l'\omterm{def:ml:feature-engineering-selection-and-importance:substitution}{effet de substitution} masque chacune derrière l'autre. Si le groupe compte, on en garde une
ou une combinaison (leur moyenne ou la première composante principale), choisie pour sa stabilité et son coût, non pour son
importance isolée.
\iqlookfor{l'\omterm{def:ml:feature-engineering-selection-and-importance:substitution}{effet de substitution} et l'importance par grappes.}
\end{solution}

\begin{solution}{iq:ml:feature-engineering-selection-and-importance:3}
Les \omterm{def:ml:feature-engineering-selection-and-importance:shapley}{valeurs de Shapley} des variables pour une prédiction, la sortie moyenne du modèle conditionnée à des sous-ensembles de variables jouant le
rôle de jeu~; elles s'additionnent à la prédiction moins la moyenne. Elles induisent en erreur avec des variables corrélées (le mérite est partagé avec une copie que le modèle
utilise par hasard) et lorsqu'on les lit comme des effets causaux plutôt que comme des descriptions du modèle.
\iqlookfor{la propriété d'additivité et la mise en garde sur les variables corrélées.}
\end{solution}

\begin{solution}{iq:ml:feature-engineering-selection-and-importance:4}
En ne gardant que les variables retenues dans la plupart des demi-échantillons~; sous l'échangeabilité, le nombre attendu de variables de bruit conservées
est au plus $q^2/((2\pi_{\mathrm{thr}} - 1)p)$. Elle ne contrôle pas les oublis~: les variables dont l'effet échappe au sélecteur de base
(interactions, non-linéarités pour un lasso) ne sont jamais sélectionnées.
\iqlookfor{la borne et son angle mort.}
\end{solution}

\begin{solution}{iq:ml:feature-engineering-selection-and-importance:5}
On prend le logarithme, on la classe par date ou on la winsorise à des quantiles estimés sur les seules données d'entraînement, et l'on vérifie
que la relation de la variable transformée avec la cible reste stable~; on envisage une surprise de volume (par rapport à son propre historique) plutôt que le
niveau.
\iqlookfor{des transformations robustes ajustées sans anticipation.}
\end{solution}

\begin{solution}{iq:ml:feature-engineering-selection-and-importance:6}
Soient $x_1, x_2$ des signes indépendants $\pm1$ et $y = x_1x_2$~: chacun est non corrélé à $y$, mais ensemble ils la déterminent.
Les criblages marginaux (corrélations, tests univariés, lasso) écartent les deux~; le criblage doit admettre les paires, ou un modèle capable de
représenter les interactions doit servir de sélecteur.
\iqlookfor{l'exemple XOR et sa conséquence.}
\end{solution}
